\section{The unified taxonomy}\label{sec:taxonomy}

Harness design decomposes into 38 separately coded choices in nine layers, not into a small number of
architectures. Figure~\ref{fig:taxonomy-map} draws the structure and Table~\ref{tab:dimensions} lists
every dimension with the size of its value set. Each dimension is a question a reader can put to any
system, answer from a pinned artifact, and support with a cited line. The coded corpus fits this premise
better than a family structure. Hierarchical clustering of the coded systems falls short of both its
pre-stated floors for silhouette and bootstrap stability (verdict negative), and on pairs where both
dimensions carry a value the mean corrected Cram\'er's $V$ is 0.166 across 666 pairs
(\S\ref{sec:landscape}). Weak coupling without families is what a space of loosely related choices
looks like when it is measured.

\begin{figure}[t]
\centering
\includegraphics[width=\linewidth]{taxonomy_map}
\caption{The unified taxonomy: 38 dimensions in nine layers, generated from the frozen schema by a
script in the replication package, so the figure cannot drift from it. Top row: the layers answering the
three necessary clauses of the definition, whose composition is the step loop
$\mathbf{A}\rightarrow\mathbf{B}\rightarrow\mathbf{C}\rightarrow\mathbf{A}$. Second row: the layers
answering clauses~(iv)--(viii), which may be trivial in a given harness; a trivial value is still a
value (\texttt{none}), never \texttt{not\_reported}. \textbf{M} is detached because it carries metadata
on the system, not a clause. Each row gives a dimension's id, key and value count: $n$ for a
single-valued enumeration, $\{n\}$ for a multi-valued one, and the schema type for the four dimensions
that are not enumerations. The two dashed arrows are the deliberate clause displacements of
\textbf{F1} (termination condition, clause~(iii)) and \textbf{G4} (permission model, clause~(vi)). No
mark is carried by color alone.}
\label{fig:taxonomy-map}
\end{figure}

%%DIMTABLE%%

Table~\ref{tab:dimensions} shows an instrument built mostly from small closed lists. %%SHAPE%% Every
cell is therefore checkable against a closed list or a stated counting, dating, or pinning rule, and the table is
generated from the schema so that it cannot drift from what was coded.

A layer is the set of dimensions that answers one clause of the working definition
(\S\ref{sec:background}), and eight of the nine follow the clauses in order. Context assembly (A)
answers input assembly, tool interface (B) action exposure and execution, control loop (C)
continuation, and memory and state (D) state. Verification and repair (E) answers checking and repair,
budget and termination (F) bounding, sandbox and environment (G) isolation, and observability and
governance (H) recording. Two placements depart from that map on purpose. The termination condition
sits in F so that termination policy and budget are reported together, and the permission model sits in
G because in practice it is configured alongside isolation. The ninth layer, M, is metadata that
stratification, weighting, and the time series need. Within a layer, one dimension records one
decision, made at one identifiable place in the pinned artifact, whose value a second reader can
recover from the same evidence. Per-dimension reliability is the test of that criterion: a dimension
that bundles two decisions yields readings that agree on the evidence and disagree on the value.

That criterion is why our cuts differ from prior taxonomies. Guo et al.\ place stopping criteria under
the control loop and budget control under verification and governance~\citep{guo2026qa2task}; a coder
cannot apply that split reliably, so we separate C from F, and E from G4 and H4. Li et al.'s context
layer covers both what the model sees this step and what persists between runs, which we split into A
and D~\citep{li2026harnessengineering}. Nine of Rombaut's twelve source-code dimensions map onto ours,
and that work supplied our evidence standard of a pinned commit with a file and
line~\citep{rombaut2026insidescaffold}. The other three have no dimension of ours to map to;
\S\ref{sec:related} names them and \S\ref{sec:threats} weighs what their omission costs.

\subsection{Multi-valued dimensions}\label{subsec:multivalued}

Fifteen dimensions hold a set, because on them a harness does several things at once: it stacks
compaction mechanisms, ships several call parsers, or gates a run at more than one point. A set cell
records every value reachable through shipped configuration, with the default named in the note. Where a value set contains \texttt{none}, it is a singleton by rule and never appears
beside another value. The registered reliability measure, exact-set agreement, is harsh on these cells:
two readings that agree on all but one primitive score as a complete disagreement, so the penalty grows
with the size of the true set rather than the size of the error. Loop primitives (C1) bears the cost. It
is the one dimension below the 0.6 freeze threshold on exact-set Cohen's kappa (0.590), and it clears
the threshold per value (0.731); Supplement~S1 reports all four measures, each of them
reproducibility between two model readings rather than correctness. We keep C1 and flag it as the
weakest dimension rather than redefine it, because redefining a dimension after seeing its reliability
is the post-hoc repair this paper audits others for.

\subsection{The absence rule}\label{subsec:absence}

The taxonomy separates a design decision from a documentation gap, and the absence rule enforces the
separation. An absence value, such as \texttt{none} or \texttt{open} for network policy, is a value on a
dimension: a choice the builders made. \texttt{not\_reported} is a property of the documents: the
sources were read and are silent. The rule codes an absence value only when the evidence contains the
place where the feature would be declared if it existed and the feature is not there. The qualifying
places are enumerated: the configuration schema, the command-line flags, the tool or plugin registry,
the run loop, or a documented feature list. Without such a place the cell is \texttt{not\_reported},
and prose that never mentions a feature is never evidence that it is absent. A taxonomy that let silence
imply \texttt{none} would measure documentation while claiming to measure design, and every
distribution over its values would mix the two with no way to separate them. The rule is what lets the
under-reporting result speak about the literature rather than the systems; its history as a protocol
amendment and its measured effect on reliability are in \S\ref{sec:methods}.

\subsection{Boundary cases and known ambiguities}\label{subsec:ambiguities}

Coding two systems end to end, SWE-agent v1.1.0~\citep{yang2024sweagent} and OpenHands
v1.16.0~\citep{wang2024openhands}, raised 26 ambiguities, and each was resolved by a rule in the
coding manual rather than by changing a value set. Table~\ref{tab:ambiguities} gives all 26 with their
resolutions. Ten rows follow one of two patterns: code the shipped default and note the reachable
alternatives (rows 1 and 13), or record the closest value and write the mismatch in the note (rows 3, 5,
6, 10, 12, 16, 24, and 25). Row 17 is the one departure from default-first coding: tracing is
coded at the highest level the artifact can reach, and the note says what runs by default. Publishing
the list is part of the instrument, because a reader who disagrees with a resolution can recode that
dimension from the released evidence strings.

%%AMBTABLE%%

\subsection{Schema freeze}\label{subsec:freeze}

The schema is frozen at v1.0.0 with all 38 dimensions intact: none was added, removed, renamed, or
re-valued between the initial draft and the freeze on 2026-09-23. The freeze was gated on
per-dimension reliability rather than on a date, and the gate, Cohen's kappa of at least 0.6, was met on
37 of 38 dimensions. The one change to the cell contract came before the freeze and added
\texttt{unresolved} beside \texttt{not\_reported}, so that a cell the coder could not settle is never
pooled with a source that is silent. The replication package's schema changelog records every change
since the draft, and any later change needs a changelog line and a registration update, so that the
released dataset and this section describe one instrument. Supplement~S1 carries the
per-layer value sets, the judgment call for each layer, the decision rules, and the two worked examples.
